\documentclass[11pt]{article}
\usepackage[T1]{fontenc}
\usepackage{lmodern}
\usepackage[letterpaper,margin=0.88in]{geometry}
\usepackage{amsmath,amssymb,amsthm,mathtools}
\usepackage{booktabs,longtable,array,microtype,xcolor}
\usepackage{enumitem,listings,fancyhdr,needspace}
\usepackage[colorlinks=true,linkcolor=blue!45!black,citecolor=blue!45!black,urlcolor=blue!45!black]{hyperref}
\usepackage{bookmark}
\definecolor{ink}{RGB}{26,45,69}
\definecolor{light}{RGB}{243,246,250}
\pagestyle{fancy}\fancyhf{}
\fancyhead[L]{\small Parity norms for a ZFC consistency equation}
\fancyhead[R]{\small 23 October 2026}
\fancyfoot[C]{\thepage}
\setlength{\headheight}{14pt}
\setlength{\parskip}{4pt}
\setlength{\parindent}{0pt}
\setlist{itemsep=3pt,topsep=5pt}
\allowdisplaybreaks
\newcommand{\Z}{\mathbb Z}
\newcommand{\Q}{\mathbb Q}
\newcommand{\N}{\mathbb N}
\newcommand{\Con}{\operatorname{Con}}
\newcommand{\Pos}{\operatorname{Pos}}
\newcommand{\Nm}{\mathcal N}
\newcommand{\file}[1]{\texttt{\detokenize{#1}}}
\newtheorem{theorem}{Theorem}[section]
\newtheorem{lemma}[theorem]{Lemma}
\newtheorem{proposition}[theorem]{Proposition}
\newtheorem{corollary}[theorem]{Corollary}
\theoremstyle{remark}
\newtheorem{remark}[theorem]{Remark}
\lstset{basicstyle=\ttfamily\footnotesize,breaklines=true,columns=fullflexible,
  keepspaces=true,backgroundcolor=\color{light},frame=single,
  rulecolor=\color{black!15},showstringspaces=false}
\title{\color{ink}\bfseries Parity norms for a smaller\\Diophantine consistency equation\\[7pt]
\large Thirteen integer unknowns, degree 1,392,\\and 211 binary arithmetic operations}
\author{Research contribution prepared for Vladimir Reshetnikov\\
\small AI-assisted mathematical research for the ProveIt project}
\date{23 October 2026\\
\small Repository snapshot \texttt{cfbf118b4ad9b2c651fe96b571d9415261766328}}
\begin{document}
\maketitle
\begin{abstract}
We construct an explicit polynomial in thirteen unrestricted integer
unknowns. Its exact total degree is \(1392\), and an attached straight-line
program evaluates it in \(211\) binary additions, subtractions, and
multiplications. Its absence of integer zeros is equivalent to
\(\Con(\mathrm{ZFC})\), by a new arithmetic reduction composed with the
retained ProveIt native/Pell/refutation calibration.

The main device combines five square conditions with fixed powers of two.
Odd squareclass coefficients cannot cancel when their powers of two are
distinct. This replaces large real separating weights by a parity argument.
The main Pell congruence is encoded in the center of the same norm, removing
its separate witness. A polynomial change of the remaining norm coordinate
then lowers degree from \(1888\) to \(1392\).

The contribution is an ordinary proof, with exact algebraic and computational
certificates. It does not supply a new end-to-end proof-assistant certificate
for the inherited native compiler. It neither assumes ZFC consistency nor
asserts global optimality or worldwide priority.
\end{abstract}

\section{Result, comparison, and evidence boundary}

There is no single order called \emph{simplest} on Diophantine equations.
Unknown count, degree, arithmetic circuit size, coefficient data, and proof
complexity are separate measures. The principal new point is
\[
  \boxed{13\Z\text{-unknowns}\quad/\quad \deg 1392\quad/\quad
          211\text{ binary operations}.}
\]
The unshifted version costs \(205\) operations at degree \(1888\). Both
polynomials use the same six literal fixed integers. Their coefficient
payloads total \(615{,}883\) uncompressed bytes; they are not hidden
existential parameters.

\begin{table}[ht]
\centering\small
\begin{tabular}{lrrr}
\toprule
Construction & Integer unknowns & Exact degree & Operations\\
\midrule
Placed source 67, Jones reduction &13&10,461,744&519\\
New five-root norm, unshifted &13&1,888&205\\
\textbf{New five-root norm, shifted}&\textbf{13}&\textbf{1,392}&\textbf{211}\\
New four-root companion &14&862&150\\
New three-root companion &15&204&128\\
New centered two-root companion &16&120&123\\
Incoming Pell-congruence parent &17&88&not used here\\
\bottomrule
\end{tabular}
\caption{Like-domain comparisons. The new companions trade more unknowns
for lower degree. Counts charge every live binary \(+,-,\times\), including
squares and multiplication by fixed numerals. Reuse and fixed leaves are free.}
\end{table}

The main degree improves the placed thirteen-integer comparison by a factor
of \(10{,}461{,}744/1392\), about \(7516\). Its operation count falls from
\(519\) to \(211\), about \(59.3\%\). These statements concern these measured
coordinates, not coefficient size or every construction in the repository.
In particular, nine-positive and ten-integer equations remain available.
A ten-positive incoming equation also exists. We do not claim a new record
for the absolute number of unknowns.

The repository was inspected at the snapshot printed on the title page.
The source directories and all eleven incoming entries were inventoried.
The new parent is the seventeen-integer degree-88 packet in
\file{con_zfc_pell_congruences_2026-10-23.zip}, including its coherent
source-34 coefficient tuple \cite{repo,parent,source67}. The archive contains
the complete parent package, its proof dependencies, actual fixed numerals,
and independent audits. New formulas are given in this report and in literal
JSON arithmetic programs.

\paragraph{Logical status.}
The new square-combination, denominator, coordinate, and degree arguments
below are proved for unbounded integer inputs. The native interpretation of
the inherited source is imported with its ordinary mathematical proof.
The resulting theorem is intended as a derivation formalizable in ZFC.
The existing project's incomplete internal-model/compiler formalization is
not completed by finite tests or by this report. This distinction matters:
an ordinary proof composed with an existing theorem, a checked arithmetic
identity, and a checked encoded ZFC derivation are different claims.

\section{The exact thirteen-variable polynomial}

\subsection{Variables and constants}

There are exactly thirteen independently quantified variables:
\[
 j,\ C,\ z,\ h,\ r,\ v,\ w,\ T,\ c_0,\ s,\ p_1,\ p_2,\ t\in\Z .
\]
Every other symbol in this section is a polynomial abbreviation.
In the machine-readable files the names \(C,v,w,s,t\) appear as
\file{C_native}, \file{v_scale}, \file{w3}, \file{pair_root},
and \file{t_parity}.

There are six fixed positive integer coefficients
\[
 \beta,\quad K_0,\quad M_C,\quad M_F,\quad A_F,\quad A_C.
\]
They are the unchanged source-34 tuple and the source-71 determinant scales:
\[
 \beta=2^{1064921}-1,\qquad A_F=2^{1113},\qquad A_C=2^{666104}.
\]
The other three are fixed finite sums over the supplied literal K215 tile
table, specified by the attached source-34 builder. Section~\ref{sec:coeff}
explains their complete binding. In particular \(M_F\) is already shifted:
\(M_F=M_F^{\rm native}+\beta\). It must not be shifted a second time.

\Needspace{14\baselineskip}
\subsection{Polynomial registers}

Define, in the displayed order,
\begin{align}
 J&=6j+1,&q&=\beta J+1,&W&=4(3w+1),\nonumber\\
 c&=2c_0+1,& i&=2r+1,&
 F&=(K_0+W)C-z(q-1),                                      \label{eq:register1}\\
 X&=qW,&Y&=(2v+1)q^3,&E&=XY,\nonumber\\
 a&=E+Y,&H&=4a+3,&\Delta&=a^2+H=(a+2)^2-1,\nonumber\\
 R&=\bigl(q(q-F)-C\bigr)(q^2-1)+(M_C+qM_F)J,\nonumber\\
 k&=R+1+hE,&\eta&=c-kY,&\zeta&=k-\eta.                      \label{eq:register2}
\end{align}
The normalized strong and auxiliary registers are
\begin{align}
 F_2&=1+\Delta i^2c^4,& I&=\Delta i c^2,\nonumber\\
 V&=F_2(cT-R)-c,&Q&=I^2,&\delta&=Q-1,\nonumber\\
 r_a&=QV^2-1,& C_a&=\delta r_a.                             \label{eq:register3}
\end{align}
The square radicands and the common-sign guard are
\begin{align}
 A_0&=1+XY^2(XY^2+1)k^2,\nonumber\\
 G&=(A_FF)(A_CC)-(q-A_FF-A_CC)^2,&Z_\eta&=\eta\zeta,\nonumber\\
 b_g&=GZ_\eta-(s^2-G-Z_\eta+1)^2-1-p_1^2-p_2^2,\nonumber\\
 B_g&=4b_g+1,&D_0&=1+\Delta c^2,&b_0&=X+ac.                 \label{eq:register4}
\end{align}
Finally set
\begin{align}
 Z_0&=H(2t+1)+b_0,&
 L&=8H(2Q-1)V,&Z&=Z_0+L,\nonumber\\
 d_0&=H^2A_0,&d_1&=D_0,&d_2&=16H^2B_g,\nonumber\\
 d_3&=16H^2F_2,&d_4&=256H^2C_a.                            \label{eq:register5}
\end{align}

\subsection{A polynomial-only norm definition}

For fixed \(m\), let
\[
 \Nm_m(Z;d_0,\ldots,d_{m-1})
   =\prod_{\varepsilon\in\{-1,1\}^m}
     \left(Z+\sum_{\ell=0}^{m-1}\varepsilon_\ell\sqrt{d_\ell}\right).
                                                               \label{eq:formalnorm}
\]
This displayed product explains the proof, but it is an ordinary polynomial
with integer coefficients: changing any formal square-root sign permutes its
factors, so only even powers of every formal root remain. Replacing each
root square by \(d_\ell\) gives the polynomial. All signs and all limits are
fixed finite sets.

\begin{minipage}{\linewidth}
For a definition without any radical notation at all, use the following
four-root evaluator. For inputs \(x,a,b,c,d\), put
\begin{align*}
 p&=ab,&q_*&=cd,&u&=a+b,&v_*&=c+d,\\
 \sigma&=x^2-u-v_*,&
 K_*&=\sigma^2-4uv_*+4(p+q_*),\\
 L_*&=-4(\sigma+2v_*),&
 M_*&=-4(\sigma+2u),&N_*&=-8,\\
 E_*&=K_*^2+pL_*^2-q_*(M_*^2+pN_*^2),\\
 O_*&=2K_*L_*-2q_*M_*N_* .
\end{align*}
Then
\[
 \Nm_4(x;a,b,c,d)=E_*^2-pO_*^2.                              \tag{N4}
\]
\end{minipage}
\medskip

To obtain \(\Nm_5\), perform these same additions and multiplications on
ordered pairs with the rules
\[
 (A,B)+(C,D)=(A+C,B+D),\qquad
 (A,B)(C,D)=(AC+eBD,\,AD+BC).
\]
Evaluate (N4) at \(x=(Z,1)\), the four radicands
\((d_0,0),\ldots,(d_3,0)\), and \(e=d_4\). If its result is the pair
\((U_*,V_*)\), define
\[
 \boxed{\Nm_5(Z;d_0,\ldots,d_4)=U_*^2-d_4V_*^2.}              \tag{N5}
\]
These are finite polynomial definitions over \(\Z\). The pairs are
arithmetic registers, not Diophantine unknowns. The evaluator never computes
a square root, a quotient, or a variable exponent.

The proposed equation is simply
\begin{equation}
 \boxed{P_{13}(j,C,z,h,r,v,w,T,c_0,s,p_1,p_2,t)
       :=\Nm_5(Z;d_0,d_1,d_2,d_3,d_4)=0.}                   \label{eq:main}
\end{equation}
There is no additional equation, positivity requirement, or hidden quantified
coordinate. The unshifted companion replaces \(Z\) by \(Z_0\).

\section{Why parity makes the norm faithful}

A norm equation alone generally does not force its constituent radicands
to be squares. For example \(\sqrt5-\sqrt5+1=1\); an unweighted
three-root norm can have an integer root even when two of its radicands are
nonsquares. The following restricted lemma is the new useful mechanism.

\begin{lemma}[Parity-separated squareclasses]\label{lem:parity}
Let \(m\ge 2\), let \(a_0,\ldots,a_{m-2}\) be odd integers, and let \(a_{m-1}\) be an arbitrary
integer. A rational value of a signed sum
\[
 \sum_{\ell=0}^{m-1}2^\ell\varepsilon_\ell\sqrt{a_\ell},
 \qquad \varepsilon_\ell\in\{-1,1\},
\]
is possible only if every \(a_\ell\) is an integer square.
Conversely, if they are all squares, its signed sums are integers, and all
are odd.
\end{lemma}

\begin{proof}
Each nonzero integer radicand has a unique expression \(d b^2\), with
\(d\) a signed squarefree integer and \(b\) a positive integer.
For an odd radicand, both \(d\) and \(b\) are odd.
Group the signed radicals by \(d\). Distinct signed squareclasses are
linearly independent over \(\Q\) in the multiquadratic field they generate.
This follows, for example, by projecting onto the characters of its finite
Galois group. Hence the coefficient of every nontrivial squareclass must
vanish in a rational sum.

But in any class containing one of the first \(m-1\) terms, its term with
the smallest power of two has an odd coefficient after dividing by that
power. All later terms have even coefficients. Their signed sum cannot
vanish. A class containing only the last nonzero term cannot vanish either.
Zero last radicands can be omitted. Therefore no nontrivial squareclass
occurs. Negative radicands would have a nontrivial signed class too, so
they are also excluded. Rational square roots of integers are integers.
The first root is odd and all subsequent weighted roots are even, proving
the converse parity statement.
\end{proof}

No order or positivity of the radicands is assumed in this lemma.
It remains valid before the guard is known to be positive.
The weights \(1,2,4,\ldots\) are constants, independent of the sizes of
the radicands.

For the actual five radicands,
\[
 d_0=H^2A_0,\quad
 d_1=4(D_0/4),\quad
 d_2=16H^2B_g,\quad
 d_3=64(H^2F_2/4),\quad
 d_4=256H^2C_a.
\]
The first four unweighted radicands are odd integers, as proved next.
The divisions by four in this sentence are proof notation only; the
literal polynomial uses \(D_0\) and \(16H^2F_2\) directly.

\section{Unconditional arithmetic facts and denominator recovery}

\begin{lemma}[Ambient congruences]\label{lem:ambient}
For every assignment of the thirteen integer inputs,
\[
 q\equiv0\pmod2,\quad X\equiv Y\equiv a\equiv0\pmod8,\quad
 \Delta\equiv3\pmod8,\quad H\text{ odd},
\]
\[
 A_0\text{ odd},\qquad B_g\equiv1\pmod4,\qquad
 D_0\equiv F_2\equiv4\pmod8.
\]
Moreover \(H\ne0\), \(\Delta>1\), \(I\ne0\), and \(Q>1\).
\end{lemma}
\begin{proof}
The coefficient \(\beta\) and \(J=6j+1\) are odd, so \(q\) is even.
The nonzero odd integer \(J\) has absolute value at least one, and the
fixed \(\beta\) is very large; in particular \(q\ne0\) and \(|q|>4\).
The integer \(W=4(3w+1)\) is a nonzero multiple of four.
Thus \(X=qW\) is a multiple of eight, and the odd multiple
\(Y=(2v+1)q^3\) is a nonzero multiple of eight.
This gives \(a=Y(X+1)\equiv0\pmod8\).
Since \(|Y|\ge|q|^3\) and \(X+1\ne0\), \(|a|\ge|q|^3>4\);
hence \((a+2)^2-1>1\).
Now \(c,i\) are odd and nonzero, so \(I=\Delta i c^2\ne0\) and
\(|I|>1\).
The remaining congruences follow by reducing the displayed definitions
modulo eight, using \(c^2\equiv i^2\equiv c^4\equiv1\pmod8\).
\end{proof}

The rational auxiliary radicand \(r_a/\delta\) is multiplied by \(\delta^2\), giving \(C_a=\delta r_a\).
Square recovery gives a rational auxiliary root; the following elementary
lemma proves that it is integral.

\begin{lemma}[Rational adjacent-Pell descent]\label{lem:descent}
Let \(Q\ge2\) and \(V\in\Z\). If \(y\in\Q\) satisfies
\[
 QV^2-(Q-1)y^2=1,
\]
then \(y\in\Z\).
\end{lemma}
\begin{proof}
Write \(y=a/b\) in lowest terms, \(b>0\).
Then \(b^2\mid Q-1\), so \((Q-1)y\) is an integer.
Take \(V>0,y\ge0\) by signs. The equation excludes \(V=0\), and
\(V=1\) gives \(y=1\).
For \(V\ge2\), \(y>V\). Define
\[
 V'=(2Q-1)V-2(Q-1)y,\qquad
 y'=(2Q-1)y-2QV.
\]
The first quantity is integral. Both satisfy the same norm equation.
Also \(V'<V\), since \(y>V\). To see \(V'>0\), observe
\[
 (2Q-1)^2V^2-4(Q-1)^2y^2=V^2+4(Q-1)>0.
\]
Consequently repeated replacement by \((V',|y'|)\) strictly decreases
the positive integer \(V\) until it reaches one.
The transformation has an inverse with integer coefficients.
At the terminal pair \(y=1\); reversing every step proves that the
original rational \(y\) was integral.
\end{proof}

\begin{remark}
No squarefreeness of \(Q-1\) is used. Merely clearing a denominator
without this argument would leave a proof gap.
\end{remark}

\section{The imported source and the new reconstruction}

\subsection{The parent interface}

The retained seventeen-integer source uses the same \(\beta,K_0,M_C,M_F\),
the same determinant scales, and the same polynomial transport
\[
 F=(K_0+W)C-z(q-1).
\]
Its free main root is represented by \(D=X+ac+H\gamma\).
It has odd scale parameters in \(Y\) and in an ordinary strong coordinate
\(I_{\rm old}=(2r_{\rm old}+1)c^2\).
Besides its monic registers and transport, its relevant equations can be
written
\begin{align}
 \tau^2&=A_0,&D^2&=1+\Delta c^2,\nonumber\\
 I_{\rm old}^2&=\Delta(f^2-1),&
 I_{\rm old}^2V^2-(I_{\rm old}^2-1)y^2&=1,                  \label{eq:parentnorms}\\
 V&=fcT_{\rm old}-f^2R-c,\nonumber\\
 GZ_\eta-(s^2-G-Z_\eta+1)^2
    &=1+p_1^2+p_2^2+p_3(p_3+1).                           \label{eq:parentguard}
\end{align}
The packet centers its first and auxiliary roots; the inverse changes are
\[
 \tau_{\rm packet}=2\tau-2P+k,\qquad
 y_{\rm gap}=y-V,\qquad P=k+XY^2k.
\]
These are integral polynomial expressions in the restored roots.

The parent report proves that its integer zeros correspond to finite
refutations of the specified ZFC calculus. Its main ingredients are
the nearest-square joint-positivity guard, determinant sign recovery,
signed Pell rank recovery, residue orientation of \(X\), a CRT
compression of two auxiliary congruences, and the fixed native-window
compiler \cite{parent,source34,source70,source71}.
The new construction restores this exact interface rather than silently
switching to another coefficient tuple.

\subsection{Every new zero restores an integer parent witness}

\begin{theorem}[Soundness of the new arithmetic reduction]\label{thm:sound}
Every integer zero of~\eqref{eq:main} gives an integer zero of the retained
seventeen-integer source.
\end{theorem}
\begin{proof}
Let \(x=2t+1+8(2Q-1)V\), so \(Z=Hx+b_0\) and \(x\) is odd.
The norm factorization and Lemmas~\ref{lem:parity}--\ref{lem:ambient}
force all five radicands to be integer squares.
Dividing their integer roots by the displayed square factors, using
rational-square integrality, gives signed integers
\[
 \tau,\ D,\ d,\ f,\ z_a
\]
such that
\[
 \tau^2=A_0,\quad D^2=D_0,\quad d^2=B_g,\quad
 f^2=F_2,\quad z_a^2=C_a.
\]
Signs can be chosen to make the vanishing factor read
\begin{equation}
 Hx+b_0=H\tau+D+4Hd+4Hf+16Hz_a.                            \label{eq:rootsum}
\end{equation}
In particular
\[
 \gamma=\frac{D-b_0}{H}=x-\tau-4d-4f-16z_a
\]
is an integer. Thus the actual main congruence is restored, not just
the main square test. The root \(d\) is odd, so
\(p_3=(d-1)/2\) is an integer. Its square identity gives exactly
\eqref{eq:parentguard}.

Set \(y=z_a/\delta\). Since \(z_a^2=\delta(QV^2-1)\),
\[
 QV^2-(Q-1)y^2=1.
\]
Lemma~\ref{lem:descent} makes \(y\) integral.
For the ordinary strong source put
\[
 r_{\rm old}=\frac{\Delta i-1}{2},\qquad T_{\rm old}=fT.
\]
These are integers, since \(\Delta i\) is odd. Then
\[
 (2r_{\rm old}+1)c^2=I,\qquad
 I^2=\Delta(F_2-1)=\Delta(f^2-1),
\]
and
\[
 fcT_{\rm old}-f^2R-c=f^2(cT-R)-c=V.
\]
The old affine \(J\)-coordinate is \(2j\); its old \(c\) coordinate is
\(2c_0+1\). The centered-root inverses above complete all seventeen
coordinates. Every parent residual is now zero.

For completeness of the sign discussion, \eqref{eq:parentguard} has a
strictly positive right side. The inherited nearest-square guard therefore
gives \(G>0\) and \(\eta\zeta>0\), and the determinant guard recovers the
required common sign of \(F,C,q-F-C\). Those conclusions are obtained
before invoking the source's signed rank and native semantics.
\end{proof}

The proof never assumes that arbitrary inputs are positive, that the
radicands are positive before extraction, or that a canceled denominator
is nonzero without a reason.

\section{Completeness without an extra congruence promise}

The normalized chart is a genuine restriction of the set of parent
witnesses. Thus a proof that every parent zero has a polynomial inverse
would be incorrect. What is needed is that every hypothetical finite
refutation supplies at least one canonical witness in this chart.
We prove that statement explicitly.

The retained source-34 canonical family has
\[
 q=2^\lambda,\quad \lambda=1064921N\text{ odd},\quad
 R\equiv3\pmod4,\quad R>3q,
\]
\[
 X=2^R,\quad Y/q^3\text{ odd},\quad
 A=a+2\text{ even},\quad c=\psi_A(R)\text{ odd}.
\]
Here \(\chi_A(n)+\psi_A(n)\sqrt{A^2-1}
=(A+\sqrt{A^2-1})^n\).
It also supplies both the determinant guard and the source residue
conditions simultaneously. In particular \(J\equiv1\pmod6\),
and \(W=4(3w+1)\).
These are documented properties of the actual same source-34 tuple,
not a borrowed property of the smaller compact tuple.

\begin{lemma}[An odd normalized strong witness]\label{lem:strongcomplete}
For this canonical \(c=\psi_A(R)\), there are an odd positive integer \(i\)
and an integer \(f>0\) satisfying \(f^2-\Delta i^2c^4=1\).
\end{lemma}
\begin{proof}
Choose \(m=Rc\), \(f=\chi_A(m)\), and \(z_*=\psi_A(m)\).
Let \(d=\chi_A(R)\). The composition identity gives
\[
 z_*=c\,\psi_d(c).
\]
Since \(d^2-1=\Delta c^2\), the binomial formula for the Pell sequence
gives
\[
 \psi_d(c)\equiv c\,d^{\,c-1}\pmod{c^2}.
\]
Consequently \(c^2\mid z_*\).
As \(A\) is even and \(m=Rc\) is odd, \(z_*\) is odd.
Thus \(i=z_*/c^2\) is a positive odd integer, and the Pell identity
is the required normalized strong equation.
\end{proof}

\begin{lemma}[The squared-modulus auxiliary choice]\label{lem:auxcomplete}
With the preceding canonical witnesses, put \(S=\Delta i c^2\).
There are positive integers \(T,V,y\) such that
\[
 V=f^2(cT-R)-c,\qquad S^2V^2-(S^2-1)y^2=1.
\]
\end{lemma}
\begin{proof}
Write \(R=2v_*+1\), where \(v_*\) is odd, and define
\[
 V_*=\chi_S(R)/S,\qquad y_*=\psi_S(R).
\]
The first expression is an integer polynomial \(Q_{v_*}(S^2)\), because
\(\chi_S(R)\) is an odd polynomial in \(S\). The Pell identity supplies
the stated auxiliary equation. The elementary odd-index polynomial
identities are
\[
 Q_{v_*}(0)=(-1)^{v_*}R=-R,\qquad
 Q_{v_*}(1-A^2)=(-1)^{v_*}\psi_A(R)=-c.
\]
They follow by the binomial formulas for \(\chi\) and \(\psi\), or by
their recurrences. Since
\(S^2=\Delta(f^2-1)\equiv1-A^2\pmod{f^2}\) and \(c\mid S\),
\[
 V_*\equiv-c\pmod{f^2},\qquad V_*\equiv-R\pmod c.
\]
Also \(f^2\equiv1\pmod c\), hence \(\gcd(c,f^2)=1\).
Therefore
\[
 T=\frac{V_*+c+Rf^2}{cf^2}
\]
is a positive integer. It yields the asserted \(V=V_*\) and \(y=y_*\).
This is precisely the squared-modulus reconstruction proved in the
retained source-61 lemma \cite{source61}; it does not demand an
unproved choice of arbitrary congruences.
\end{proof}

\begin{theorem}[Completeness of the new arithmetic reduction]\label{thm:complete}
A finite refutation of the specified ZFC calculus yields a thirteen-integer
zero of~\eqref{eq:main}.
\end{theorem}
\begin{proof}
Take the simultaneous canonical native witnesses.
Set \(j=(J-1)/6\), \(c_0=(c-1)/2\), and
\(r=(i-1)/2\), using Lemma~\ref{lem:strongcomplete}.
Reselect the auxiliary witnesses by Lemma~\ref{lem:auxcomplete}.
Keep the native transport and guard coordinates.
The ternary positive representation provides \(p_1,p_2,p_3\), and put
\(d=2p_3+1\).
Let \(\tau,D,\gamma\) be the canonical first and main roots and the main
quotient.

The congruences give \(\tau\) odd, \(D\) even, and \(b_0=X+ac\) even.
Since \(H\) is odd, \(\gamma=(D-b_0)/H\) is even.
Thus
\[
 x=\tau+\gamma+4d+4f+16\delta y
\]
is odd. Define
\[
 t=\frac{x-1}{2}-4(2Q-1)V\in\Z.
\]
Then \eqref{eq:rootsum} holds with these signs and
\(z_a=\delta y\). A factor of the norm vanishes. This constructs an
actual integer solution from the hypothetical finite certificate.
\end{proof}

\begin{corollary}[Consistency equivalence]\label{cor:con}
Composing Theorems~\ref{thm:sound} and~\ref{thm:complete} with the retained
native/Pell/refutation theorem gives, by ordinary mathematics formalizable
in ZFC,
\[
 \mathrm{ZFC}\vdash
 \left(\neg\exists\mathbf x\in\Z^{13}\;P_{13}(\mathbf x)=0\right)
 \longleftrightarrow \Con(\mathrm{ZFC}).
\]
\end{corollary}
There is no appeal to the truth of either consistency or inconsistency.
All completeness choices are conditional constructions from a finite
refutation. The theorem does not conflict with the second incompleteness
theorem: it proves an equivalence, not either side of it.

\section{The degree reduction and exact degree certificate}

The unshifted norm uses \(Z_0\). The shifted version is obtained by the
integer polynomial change
\[
 t_{\rm old}=t+4(2Q-1)V.
\]
All registers on the right are independent of \(t\), so this is a triangular
polynomial bijection with inverse subtraction. Therefore the two
thirteen-variable polynomials have equivalent zero sets, with no added
number-theoretic condition.

With inputs of degree one and fixed coefficients of degree zero,
\[
\begin{array}{c|rrrrrrrrrrrr}
\text{register}&q&X&Y&E&a&\Delta&R&k&I&F_2&V&C_a\\ \hline
\deg&1&2&4&6&6&12&5&7&15&18&23&106
\end{array}
\]
The degrees of \(G,Z_\eta,B_g\) are \(4,22,44\), respectively.
Thus the radicand degrees are
\[
 (\deg d_0,\ldots,\deg d_4)=(46,14,56,30,118).
\]
The carrier \(Z_0\) has degree seven. In the unshifted product, the last
radicand strictly dominates; its sixteenth power gives exact degree
\(16\cdot118=1888\).

For the shifted norm, pair the two signs of its final radical first.
For a signed sum \(S_4\) of the other four radicals, the paired factor is
\[
 (Z_0+L-S_4)^2-d_4.
\]
The critical identity is
\begin{equation}
 L^2-d_4=64H^2V^2+256H^2\delta.                            \label{eq:midpoint}
\end{equation}
It is an exact polynomial identity, since
\((2Q-1)^2-4Q(Q-1)=1\). Therefore the paired factor equals
\[
 (Z_0-S_4)^2+2L(Z_0-S_4)+64H^2V^2+256H^2\delta.
\]
Use weights one-half of each radicand degree for its formal root.
The largest root among the four is \(\sqrt{d_2}\), of weight \(28\);
\(\deg L=59\). Every paired factor has weight at most \(87\), with
unique highest term \(-2L(\pm\sqrt{d_2})\).
There are sixteen paired factors. Their product has degree at most
\(16\cdot87=1392\), and its highest homogeneous part is
\[
 \boxed{2^{16}[L]_{59}^{\,16}[d_2]_{56}^{\,8}.}
\]
This is nonzero for the actual fixed coefficient tuple.
The attached independent certificate substitutes specified positive integer
slopes times an indeterminate into every input and reduces the actual six
coefficient numerals modulo \(1{,}000{,}000{,}007\). The degree-1392
coefficient is
\[
 756{,}708{,}762\ne0\pmod{1{,}000{,}000{,}007}.
\]
Together with the proved upper bound, this is an exact nonvanishing
certificate, not a random estimate of degree. A separate finite-difference
calculation reproduces it.

\section{The smaller companion equations}

The archive retains intermediate constructions because a higher unknown
count can buy a much lower degree.

\subsection{A two-root merger in sixteen integers}

In the seventeen-variable parent, \(\tau^2=A=4a_4+1\) and
\(\delta y^2=r_a\), with \(\delta\ne0\).
For \(x=2t+1=\tau+2y\), clearing the two-root norm and dividing its
fixed content gives
\[
 N_{16}=\bigl[\delta(t(t+1)-a_4)-r_a\bigr]^2-A\delta r_a.
\]
A zero gives
\[
 \tau=\frac{\delta(x^2+A)-4r_a}{2x\delta}\in\Q,\qquad
 \tau^2=A.
\]
Thus \(\tau\) is an odd integer and \(y=(x-\tau)/2\) is an integer.
This particular two-root version needs no adjacent-Pell descent.

For the centered form put \(z_0=XY^2/4\), a polynomial because
\(X=4q(3w+1)\), and substitute the even center \(4z_0k\) into \(x\).
The resulting numerator is
\[
 \delta\bigl[t(t+1)+4z_0kt+2z_0k-z_0k^2\bigr]-r_a.
\]
Combining its norm with the three retained parent rows by a sum of squares
gives degree \(120\) in sixteen integers, evaluated in \(123\) operations.

\subsection{Three and four roots}

The fifteen-integer companion also packs the guard root
\(2p_3+1\) with weights \(1,2,4\).
The three residual degrees are \(14,14,176\).
For any integer residuals \(r_1,\ldots,r_n\),
\[
 \prod_{\ell=1}^n(1+4r_\ell)-1=0
 \quad\Longleftrightarrow\quad r_1=\cdots=r_n=0,             \label{eq:unit}
\]
because every factor must be an integer unit, and no factor can equal
\(-1\) modulo four. This yields degree \(204\) and \(128\) operations.

The fourteen-integer companion uses the normalized strong chart, combines
four roots, and retains the main congruence witness \(\gamma\).
Its two residual degrees are \(14,848\); \eqref{eq:unit} gives degree
\(862\) and \(150\) operations. The thirteen-integer construction then
absorbs \(\gamma\) into the center and needs no residual finalizer.

These intermediate results explain how the main construction was obtained.
They do not make the final polynomial depend on an optimization search.

\section{Fixed coefficient specification and provenance}\label{sec:coeff}

The six values are genuine integer coefficients fixed before any question
about halting is asked.
Their storage and reproducibility are part of the mathematical artifact.

\begin{center}
\begin{tabular}{lrr}
\toprule
Coefficient&Bits&Uncompressed bytes\\\midrule
\(\beta\)&1,064,921&133,116\\
\(K_0\)&1,065,043&133,131\\
\(M_C\)&1,064,921&133,116\\
\(M_F\)&1,064,922&133,116\\
\(A_F\)&1,114&140\\
\(A_C\)&666,105&83,264\\\midrule
Total&&615,883\\\bottomrule
\end{tabular}
\end{center}

The files under \file{coefficients/payloads/} are minimal unsigned
little-endian representations, gzip-compressed. The binding manifest gives
both raw and compressed SHA-256 hashes. The source-34 recipe is included,
with its literal \file{k215.tm} and \file{wang943_tiles.json}. The recipe
checks all 430 transitions and all 943 generated tiles before constructing
the coefficient sums.

For a compact mathematical summary, it uses radix
\(\rho=2^{121}\), block length \(8801\), and \(B=\rho^{8801}\).
It builds the fixed selector weights, low mask, high mask, and convolution
polynomial \(D_C\) from the authenticated tile list, with \(D_R=\rho\), then
\[
 \beta=B-1,\qquad K_0=D_C+BD_R,\qquad
 M_F=M_F^{\rm native}+B-1.
\]
The two determinant scales are the powers of two already printed above.
Every sum and loop in the builder has a fixed finite range; it never waits
for K215 to halt. Its exact literal transition hash is
\begin{center}\scriptsize\ttfamily
fdbcef0825886035c3091c8537a69535bc96f9e265f0ae45863543d40ae70090
\end{center}

The parent packet is retained byte for byte.
No coefficient from the compact source-44 or original source-31 tuple is
substituted for a source-34 port. Equality of radix exponents would not
justify such a substitution.

The archive includes the complete incoming parent package as
\file{upstream/pell-congruences-source/}, with its original notices and
proofs. This retains the exact source semantics and canonical witness
promises required above. A machine-readable source ledger pins
the Git commit, archive hash, parent packet, and coefficient binding.

\section{Verification and how to reproduce the result}

The checks address concrete failure modes.

\begin{enumerate}
\item \textbf{Universal norm identities.}
An independent generator expands the generic four- and five-root
polynomials by Newton identities. Their \(495\) and \(20{,}349\) monomials
agree exactly with the independently built multiquadratic norm circuits.
The compact four-root evaluator is also checked by its biquadratic
coefficient identity; lifting it to the final quadratic algebra proves
the five-root evaluator.
\item \textbf{Literal graph and coefficient checks.}
The delivered program contains only live binary \(+,-,\times\) nodes.
It has \(211\) operations: \(51\) additions, \(31\) subtractions,
and \(129\) multiplications. All fixed coefficients are authenticated.
Independent graph evaluation agrees with the separate source packets.
The count is an attained upper bound, not a claim that no shorter circuit
exists.
\item \textbf{Exact degree.}
Safe degree bounds, the cancellation-resolving identity
\eqref{eq:midpoint}, a nonzero specialized leading coefficient, and
independent finite differences agree on \(1392\).
\item \textbf{Exceptional cases and negative controls.}
Tests include signed radicands, zero guard candidates, cleared-denominator
degeneracies, omitted oddness, and the false unweighted norm.
The proof excludes each relevant false branch for unbounded inputs.
The parent two/three-root companions explicitly handle \(Q=0\) when their
broader charts allow it; the final normalized thirteen-variable chart has
\(Q>1\) on every integer assignment.
\item \textbf{Bounded diagnostics.}
The independent generic checker examines \(119{,}000\) pair assignments,
\(143{,}990\) triple assignments, \(3{,}360\) completeness witnesses, and
\(7{,}059\) rational-Pell fibers. These are transcription and adversarial
diagnostics, not evidence for ZFC consistency or substitutes for the proofs.
\end{enumerate}

From the archive root, the release runner rebuilds the equations,
the compact arithmetic programs, and the independent certificates:
\begin{lstlisting}
python3 verify_release.py
\end{lstlisting}
All new mathematics/code requires only the Python standard library.
The report builds with a conventional LaTeX installation:
\begin{lstlisting}
cd report
latexmk -pdf -interaction=nonstopmode -halt-on-error parity-norms.tex
\end{lstlisting}
The exact paths and commands are recorded in the top-level README.
The coefficient recipe and inherited release scripts keep their own
original execution conventions; they are not silently relabeled as new
proofs.

\section{What is new, and what remains open}

The useful new idea is a restricted constant-weight square-combination
theorem, applied where the source already supplies the needed parities.
The method has three parts:
\begin{itemize}
\item separate squareclass contributions by powers of two;
\item make a norm center carry a congruence as well as a sum of roots;
\item remove the leading auxiliary square by an exact midpoint shift.
\end{itemize}
The parity statement is proved here and appears absent from the inspected
ProveIt source and incoming constructions in this role.
We do not claim that the general algebra of multiquadratic norms is new,
or that this is the first use of parity in Diophantine representation.
Literature searches were not a proof of historical priority.

The construction retains K215 as the inherited source of one concrete
coefficient tuple and its refutation calibration. Its new compression
does not use Jones's nine-variable reduction or Sun's ten-integer
reduction. The square-combination lemma is independent of that source and
may be reusable for other arithmetic encodings.
Sun's ten-integer theorem and the universal-pair literature remain the
appropriate broader comparisons \cite{sun,bayer}.

The principal open tasks are precise:
\begin{enumerate}
\item Formalize the parity lemma, normalized chart, reconstruction, and
midpoint identity inside the existing proof development.
\item Close the inherited native/compiler/internal-model calibration if an
end-to-end checked ZFC derivation is desired.
\item Investigate whether another guarded congruence can be moved into the
same center. The transport modulus \(q-1\) need not be coprime to \(H\);
canceling it without a proof would be invalid.
\item Reduce the guard's degree or its witnesses. A claimed two-variable
replacement for the ternary positive representation needs its own
surjectivity theorem.
\end{enumerate}
Neither absolute minimality nor a smaller-variable universal representation
is established by the present result. The delivered thirteen-variable
polynomial is a concrete, reproducible improvement at a different point
of the repository's size--degree frontier.

\begin{thebibliography}{99}
\raggedright
\bibitem{repo}
V.\ Reshetnikov et al., \emph{ProveIt}, snapshot
\texttt{cfbf118b4ad9b2c651fe96b571d9415261766328}.
\url{https://github.com/VladimirReshetnikov/ProveIt}.
See \file{Computability/DiophantineConsistency/} and
\file{Computability/HilbertTenthProblem/}.
\bibitem{parent}
ProveIt incoming package, \emph{Pell congruences, shared guards, and monic
transport}, \file{con_zfc_pell_congruences_2026-10-23.zip}.
The complete package is retained in this release; the principal parent
is \file{transport/degree88_17.json}.
\bibitem{source34}
ProveIt source 34, \emph{A mixed-radix, parity-separated, unrounded
coefficient tuple}. The exact coefficient builder and canonical odd-scale
proof are retained in the incoming parent.
\bibitem{source61}
ProveIt source 61, \emph{Congruence carriers and a divided Pell norm},
\file{61-congruence-carriers/report/congruence_carriers.tex},
especially the normalized strong and squared-modulus auxiliary lemmas.
\bibitem{source67}
ProveIt source 67, \emph{Oriented guards},
\file{67-oriented-guards/README.md} and its thirteen-integer companion.
\bibitem{source70}
ProveIt source 70, \emph{Residue orientation}, the signed source and
removal of the separate \(X=2u^2\) witness.
\bibitem{source71}
ProveIt source 71, \emph{Determinant guards}, including canonical
completeness for the six-port source-34 tuple.
\bibitem{sun}
Z.-W.\ Sun, \emph{Ten unknowns for Hilbert's tenth problem over the
integers}, arXiv:2609.01594v2 (2026).
\url{https://arxiv.org/abs/2609.01594}.
\bibitem{bayer}
J.\ Bayer, M.\ David, M.\ Hassler, Y.\ Matiyasevich, and D.\ Schleicher,
\emph{Diophantine Equations over \(\Z\): Universal Bounds and Parallel
Formalization}, arXiv:2506.20909v3 (2026 revision).
\url{https://arxiv.org/abs/2506.20909}.
\end{thebibliography}
\end{document}

